\section*{Anhang}
\markboth{Anhang}{Anhang}

\subsection*{A\quad Domänenmodell (schema.cds)}
\begin{lstlisting}[caption={Vollständiges CDS-Schema},label={lst:anhang-schema}]
namespace purchase.approval;
using { cuid, managed, Currency } from '@sap/cds/common';

type RequisitionStatus : String(20) enum {
  Draft; Submitted; Approved; Rejected;
}

entity PurchaseRequisitions : cuid, managed {
  description       : String(256)  @mandatory;
  requester         : String(128)  @mandatory;
  costCenter        : String(10)   @mandatory;
  totalAmount       : Decimal(15,2);
  currency          : Currency;
  status            : RequisitionStatus default 'Draft';
  workflowId        : String(256);
  statusCriticality : Integer default 0 @UI.Hidden;
  items             : Composition of many Items
                        on items.parent = $self;
}

entity Items : cuid {
  parent        : Association to PurchaseRequisitions;
  materialNo    : String(40)   @mandatory;
  description   : String(256);
  quantity      : Decimal(13,3) @mandatory;
  unitPrice     : Decimal(15,2) @mandatory;
  currency      : Currency;
  netAmount     : Decimal(15,2);
}
\end{lstlisting}

\subsection*{B\quad Servicedefinition (purchase-service.cds)}
\begin{lstlisting}[caption={CDS-Servicedefinition mit Draft-Support und Aktionen},label={lst:anhang-service}]
using { purchase.approval as db } from '../db/schema';

service PurchaseService @(path: '/purchase') {
  @odata.draft.enabled
  entity PurchaseRequisitions
    as projection on db.PurchaseRequisitions actions {
      action submit() returns PurchaseRequisitions;
  };
  entity Items as projection on db.Items;

  action workflowCallback (
    workflowId : String,
    status     : String,
    comment    : String
  ) returns PurchaseRequisitions;

  function getSpaInboxUrl() returns String;
}
\end{lstlisting}

\subsection*{C\quad OpenAPI-Spezifikation (Callback-Action)}
\begin{lstlisting}[caption={OpenAPI-Spezifikation für den Workflow-Callback},label={lst:anhang-openapi}]
{
  "openapi": "3.0.0",
  "info": {
    "title": "Purchase Approval Callback",
    "version": "1.0.0"
  },
  "paths": {
    "/purchase/workflowCallback": {
      "post": {
        "operationId": "workflowCallback",
        "requestBody": {
          "content": {
            "application/json": {
              "schema": {
                "properties": {
                  "workflowId": { "type": "string" },
                  "status": {
                    "type": "string",
                    "enum": ["APPROVED", "REJECTED"]
                  },
                  "comment": { "type": "string" }
                }
              }
            }
          }
        }
      }
    }
  }
}
\end{lstlisting}

\subsection*{E\quad Submit-Aktion mit Workflow-Start (purchase-service.js -- Auszug)}
\begin{lstlisting}[caption={Submit-Aktion: Validierung, Statuswechsel und Workflow-Start},label={lst:anhang-submit}]
this.on('submit', PurchaseRequisitions, async req => {
  const id = req.params[0]?.ID || req.params[0];

  const pr = await SELECT.one.from(PurchaseRequisitions)
    .where({ ID: id })
    .columns(pr => { pr('*'); pr.items(i => i('*')); pr.attachments(a => a('*')); });

  if (!pr) return req.error(404, `Bestellanforderung ${id} nicht gefunden`);
  if (pr.status === 'Submitted')
    return req.error(409, 'Bereits eingereicht');
  if (!pr.items || pr.items.length === 0)
    return req.error(422, 'Mindestens eine Position erforderlich');
  if (pr.totalAmount <= 0)
    return req.error(422, 'Gesamtbetrag muss groesser als 0 sein');

  await UPDATE(PurchaseRequisitions).set({ status: 'Submitted' }).where({ ID: id });

  try {
    const workflowInstanceId = await workflowClient.startApprovalWorkflow(pr);
    await UPDATE(PurchaseRequisitions)
      .set({ workflowId: workflowInstanceId }).where({ ID: id });
  } catch (err) {
    await UPDATE(PurchaseRequisitions).set({ status: 'Draft' }).where({ ID: id });
    return req.error(502, `Workflow konnte nicht gestartet werden: ${err.message}`);
  }
});

this.on('workflowCallback', async req => {
  const { workflowId, status, comment, approverName, decidedAt } = req.data;

  if (!workflowId) return req.error(400, 'workflowId ist erforderlich');
  if (!['APPROVED', 'REJECTED'].includes(status))
    return req.error(400, `Ungueltiger Status "${status}"`);

  const pr = await SELECT.one.from(PurchaseRequisitions).where({ workflowId });
  if (!pr) return req.error(404, `Kein Datensatz fuer Workflow ${workflowId}`);
  if (pr.status !== 'Submitted') return req.error(409, 'Datensatz nicht im Status Submitted');

  const callbackStatus = status === 'APPROVED' ? 'Approved' : 'Rejected';
  await UPDATE(PurchaseRequisitions).set({
    status: callbackStatus,
    approverComment: comment || null,
    approvedBy: approverName || null,
    decidedAt: new Date(decidedAt || Date.now()).toISOString(),
    ...(callbackStatus === 'Rejected' ? { workflowId: null } : {})
  }).where({ ID: pr.ID });
});
\end{lstlisting}

\subsection*{F\quad WorkflowClient (workflow-client.js -- Auszug)}
\begin{lstlisting}[caption={WorkflowClient: OAuth2-Authentifizierung und Workflow-Start via SPA-API},label={lst:anhang-workflowclient}]
class WorkflowClient {

  async _getToken() {
    const response = await fetch(process.env.SPA_TOKEN_URL, {
      method: 'POST',
      headers: {
        'Content-Type': 'application/x-www-form-urlencoded',
        'Authorization': 'Basic ' + Buffer.from(
          `${process.env.SPA_CLIENT_ID}:${process.env.SPA_CLIENT_SECRET}`
        ).toString('base64')
      },
      body: 'grant_type=client_credentials'
    });
    const data = await response.json();
    return data.access_token;
  }

  async startApprovalWorkflow(requisition) {
    const token = await this._getToken();
    const workflowContext = {
      requisitionId: requisition.ID,
      description:   requisition.description,
      requester:     requisition.requester,
      costCenter:    requisition.costCenter,
      totalAmount:   parseFloat(requisition.totalAmount),
      currency:      requisition.currency_code || 'EUR',
      items: requisition.items.map(item => ({
        materialNo:  item.materialNo,
        description: item.description,
        quantity:    String(item.quantity),
        unitPrice:   String(item.unitPrice),
        netAmount:   String(item.netAmount)
      })),
      detailLink: `${process.env.APP_URL}#/PurchaseRequisitions(ID=${requisition.ID},IsActiveEntity=true)`,
      callbackUrl: process.env.CAP_CALLBACK_URL
    };

    const result = await fetch(
      `${process.env.SPA_API_URL}/workflow/rest/v1/workflow-instances`,
      {
        method: 'POST',
        headers: {
          'Authorization': `Bearer ${token}`,
          'Content-Type': 'application/json'
        },
        body: JSON.stringify({
          definitionId: process.env.SPA_WORKFLOW_DEFINITION_ID,
          context: workflowContext
        })
      }
    );
    const data = await result.json();
    return data.id;
  }
}
\end{lstlisting}

\subsection*{D\quad MTA-Deskriptor (mta.yaml -- Auszug)}
\begin{lstlisting}[caption={MTA-Deskriptor mit Modulen und Ressourcen},label={lst:anhang-mta}]
ID: purchase-approval
version: 1.0.0
modules:
  - name: purchase-approval-srv
    type: nodejs
    path: gen/srv
    requires:
      - name: purchase-approval-db
      - name: purchase-approval-auth
      - name: purchase-approval-destination
    provides:
      - name: srv-api
        properties:
          srv-url: ${default-url}

  - name: purchase-approval-db-deployer
    type: hdb
    path: gen/db
    requires:
      - name: purchase-approval-db

  - name: purchase-approval-app
    type: approuter.nodejs
    path: app
    requires:
      - name: srv-api
      - name: purchase-approval-auth

resources:
  - name: purchase-approval-db
    type: com.sap.xs.hdi-container
  - name: purchase-approval-auth
    type: org.cloudfoundry.managed-service
    parameters:
      service: xsuaa
      service-plan: application
      path: ./xs-security.json
  - name: purchase-approval-destination
    type: org.cloudfoundry.managed-service
    parameters:
      service: destination
      service-plan: lite
\end{lstlisting}
